\documentclass[paper=a4,fontsize=11pt]{jlreq}

\usepackage{amsmath}
\usepackage{hyperref}
\hypersetup{colorlinks=true,linkcolor=blue,urlcolor=blue}

\title{Japanese Book Starter}
\author{Your Name}
\date{\today}

\begin{document}
\maketitle
\tableofcontents

\section{はじめに}
これは \texttt{jlreq} 文書クラスを使った日本語の本のスターターです。
章や節を追加して、研究ノート、教科書、報告書などに利用できます。

\section{構成}
日本語の組版では、段落の間隔、禁則処理、縦組みなどを文書クラスが管理
します。本文を編集しながら、必要に応じてクラスのオプションを調整して
ください。

\subsection{数式}
数式も通常の LaTeX と同じように記述できます。
\[
  \sum_{k=1}^{n} k = \frac{n(n+1)}{2}.
\]

\section{次のステップ}
このファイルを複製して章を分けたり、参考文献や図表を追加したりできます。

\end{document}
